\section{Español abreviado}\label{sec:abreviaturas}

El contenido de los apuntes sigue en español, pero abreviado. Estas son las
abreviaturas comunes; cada persona tiene las suyas y puede añadirlas a su
leyenda personal (sección~\ref{sec:leyenda}).

\subsection{Palabras frecuentes}
\begin{multicols}{3}\small
\begin{description}[style=sameline,leftmargin=12mm,labelwidth=10mm,itemsep=1pt]
\item[q] que
\item[xq] porque
\item[x] por
\item[p/] para
\item[c/] con
\item[s/] sin
\item[tb] también
\item[cdo] cuando
\item[dnd] donde
\item[+] más
\item[\(-\)] menos
\item[d] de
\item[dl] del
\item[k] como
\item[mm] mismo
\item[aprox] aproximado
\item[máx] máximo
\item[mín] mínimo
\item[n°] número
\item[ppal] principal
\item[gral] general
\item[dif] diferencia
\item[prob] problema
\item[sol] solución
\item[info] información
\item[s.] siglo
\item[a.C.] antes de Cristo
\item[sist] sistema
\item[func] función
\item[gob] gobierno
\item[soc] sociedad
\item[econ] economía
\item[pol] política
\item[prof] profesor(a)
\item[lib] libro
\item[cap] capítulo
\end{description}
\end{multicols}

«xq» sirve para una razón o explicación dentro del discurso. Cuando se trata de
una causa en el mundo, prefiero Causa \kss{K23}.

\subsection{Terminaciones}
\begin{center}
\begin{tabular}{@{}lll@{}}
\toprule Terminación & Se escribe & Ejemplo \\ \midrule
-ción, -sión & -cn & revolucn, divisn \\
-mente & -mte & rápidamte \\
-miento & -mto & conocimto \\
-dad & -dd & socdd, velocdd \\
plural & +s & sists, funcs \\
\bottomrule
\end{tabular}
\end{center}

\subsection{Palabras largas}
En palabras largas quito vocales interiores y conservo la primera letra, las
consonantes principales y la última: «gobierno» se escribe «gbno» y «desarrollo»,
«dsrrllo». Solo lo hago si puedo leerlas sin dudar una semana después.

\subsection{Nombres del tema}
Los nombres propios o técnicos que se repiten mucho en una clase se abrevian
con iniciales y se declaran arriba de la página: «RF = Revolución Francesa»,
«MT = máquina de Turing». Esa abreviatura vale solo en esa página, salvo que la
lleve a la leyenda personal.
